\documentclass[spanish,letterpaper,oneside]{article}

% Datos del documento
\def\documenttitle {Escuela tradicional y Nueva Escuela Mexicana}
\def\documentsubtitle {Fundamentos para la enseñanza y el aprendizaje I}
\def\documentsubject {Actividad 6 · Cuadro comparativo y reflexión docente}
\def\documentauthor {Mtro. Martín Jonathan de la Cruz Muñoz}
\def\coursename {Fundamentos para la enseñanza y el aprendizaje I}
\def\coursecode {MAEM}

% Identidad institucional
\def\universityname {Instituto de Investigación, Innovación y Estudios de Posgrado para la Educación del Estado de Nuevo León}
\def\universityfaculty {Maestría en Aprendizaje y Enseñanza de las Matemáticas}
\def\universitydepartment {Fundamentos para la enseñanza y el aprendizaje I}
\def\universitydepartmentimage {img/iiiepe/marca-iiiepe.png}
\def\universitydepartmentimagecfg {height=1.75cm}
\def\universitylocation {Monterrey, Nuevo León}
\def\studentfullname {Mtro. Martín Jonathan de la Cruz Muñoz}
\def\studentid {00874}
\def\studentemail {martin.delacruz@campus.iiiepe.edu.mx}
\def\programperiod {1er tetramestre, marzo--julio de 2026}
\def\programmodality {Modalidad presencial}
\def\programshort {MAEM}

\def\authortable {
  \begin{tabular}{ll}
    \textbf{Alumno:} & \studentfullname \\
    Matrícula: & \studentid \\
    Correo institucional: & \studentemail \\
    Unidad: & \coursename \\
    Programa: & \programshort \\
    Periodo: & \programperiod \\
    & \\
    \multicolumn{2}{l}{Fecha de realización: 13 de octubre de 2026} \\
    \multicolumn{2}{l}{\universitylocation}
  \end{tabular}
}

\input{template}

\usepackage{helvet}
\renewcommand{\familydefault}{\sfdefault}
\usepackage{array}
\usepackage{longtable}
\usepackage{pdflscape}
\usepackage{tikz}
\usepackage[most]{tcolorbox}
\usetikzlibrary{arrows.meta,positioning,calc,fit,backgrounds}

\definecolor{iiiepeNavy}{HTML}{163B5C}
\definecolor{iiiepeBlue}{HTML}{1F6F9F}
\definecolor{iiiepeAqua}{HTML}{20B7A8}
\definecolor{iiiepeOrange}{HTML}{F28C28}
\definecolor{iiiepePaper}{HTML}{F6F8F7}
\definecolor{iiiepeInk}{HTML}{1E2A32}
\def\portraittitlecolor {iiiepeNavy}
\setcitestyle{authoryear,open={(},close={)}}

\newcommand{\insertcoverwatermark}{%
  \AddToShipoutPictureBG*{%
    \begin{tikzpicture}[remember picture,overlay]
      \node[opacity=.09,inner sep=0pt] at (current page.center)
      {\includegraphics[width=.58\paperwidth]{img/iiiepe/logo-puntos.png}};
    \end{tikzpicture}%
  }%
}

\newtcolorbox{reportelecturabox}[1][]{%
  enhanced,breakable,
  colback=iiiepePaper,colframe=iiiepeNavy,
  boxrule=1pt,leftrule=4pt,arc=2pt,
  left=8pt,right=8pt,top=6pt,bottom=6pt,
  before skip=6pt,after skip=10pt,
  fonttitle=\bfseries\color{white},coltitle=white,
  colbacktitle=iiiepeNavy,
  attach boxed title to top left={xshift=8pt,yshift=-2pt},
  boxed title style={colback=iiiepeNavy,colframe=iiiepeNavy,arc=1pt},
  title={Reporte de lectura},#1}

\newtcolorbox{fichalecturabox}[1][]{%
  enhanced,breakable,
  colback=iiiepeAqua!7,colframe=iiiepeAqua!70!iiiepeNavy,
  boxrule=.6pt,arc=1pt,
  left=7pt,right=7pt,top=4pt,bottom=4pt,
  fontupper=\small,before upper={\setstretch{1.0}},
  fonttitle=\bfseries\color{iiiepeNavy},coltitle=iiiepeNavy,
  colbacktitle=iiiepeAqua!7,
  title={Ficha de la lectura},#1}



\usepackage{tabularx,colortbl,pgfplots,enumitem,needspace}\pgfplotsset{compat=1.18}\definecolor{iiiepePale}{HTML}{EEF5F8}
\setlength{\parskip}{5pt}\setlength{\parindent}{0pt}
\begin{document}\insertcoverwatermark\templatePortrait\templatePagecfg\templateFinalcfg\fancyfoot{}\fancyfoot[L]{\scriptsize IIIEPE · MAEM}\fancyfoot[C]{\scriptsize\thepage}\fancyfoot[R]{\scriptsize Actividad 6}\setstretch{1.3}\raggedbottom
\section{Comparar decisiones de enseñanza y evaluación}
La comparación entre escuela tradicional y Nueva Escuela Mexicana se construye como un análisis de tendencias pedagógicas, no como una clasificación absoluta de docentes o instituciones. La lectura de evaluación incluida en EVA, basada en Scallon (2004), describe cambios hacia desempeño, autenticidad, interacción y participación del estudiante; esos cambios anteceden a la NEM y no equivalen por sí solos a todo su marco curricular. El Plan de Estudio 2022 aporta la referencia para discutir el enfoque humanista, comunitario e inclusivo.
\section{Cuadro comparativo}


{\small\setstretch{1.1}\renewcommand{\arraystretch}{1.25}
\begin{longtable}{|>{\raggedright\arraybackslash}p{\dimexpr0.3\linewidth-2\tabcolsep-2\arrayrulewidth\relax}|>{\raggedright\arraybackslash}p{\dimexpr0.35\linewidth-2\tabcolsep-2\arrayrulewidth\relax}|>{\raggedright\arraybackslash}p{\dimexpr0.35\linewidth-2\tabcolsep-2\arrayrulewidth\relax}|}\hline
\rowcolor{iiiepePale} \textbf{Criterio} & \textbf{Tendencia tradicional} & \textbf{Enfoque de la NEM} \\\hline
\endfirsthead
\hline
\rowcolor{iiiepePale} \textbf{Criterio} & \textbf{Tendencia tradicional} & \textbf{Enfoque de la NEM} \\\hline
\endhead
Finalidad & Predominio de transmisión y dominio de contenidos escolares. & Formación integral, pensamiento crítico y vinculación con la comunidad; los contenidos se articulan con problemas y contextos. \\\hline
Organización del conocimiento & Fragmentación por asignaturas y ejercicios aislados como tendencia. & Campos formativos, integración de saberes y contextualización mediante el programa analítico. \\\hline
Papel docente & Explicar, dirigir y verificar respuestas. & Diseñar situaciones, mediar, contextualizar y ajustar apoyos con autonomía profesional. \\\hline
Papel del estudiante & Recibir información y reproducir procedimientos. & Investigar, argumentar, colaborar y participar en decisiones de aprendizaje. \\\hline
\end{longtable}}



{\small\setstretch{1.1}\renewcommand{\arraystretch}{1.25}
\begin{longtable}{|>{\raggedright\arraybackslash}p{\dimexpr0.3\linewidth-2\tabcolsep-2\arrayrulewidth\relax}|>{\raggedright\arraybackslash}p{\dimexpr0.35\linewidth-2\tabcolsep-2\arrayrulewidth\relax}|>{\raggedright\arraybackslash}p{\dimexpr0.35\linewidth-2\tabcolsep-2\arrayrulewidth\relax}|}\hline
\rowcolor{iiiepePale} \textbf{Criterio} & \textbf{Tendencia tradicional} & \textbf{Enfoque de la NEM} \\\hline
\endfirsthead
\hline
\rowcolor{iiiepePale} \textbf{Criterio} & \textbf{Tendencia tradicional} & \textbf{Enfoque de la NEM} \\\hline
\endhead
Problemas matemáticos & Ejercicios con datos cerrados y una ruta anticipada. & Situaciones que permiten modelar, contrastar estrategias y justificar supuestos. \\\hline
Evaluación & Predominio de producto final, examen y calificación. & Seguimiento de procesos, diálogo, autoevaluación y uso de evidencia para retroalimentar. \\\hline
Error & Se registra principalmente como respuesta incorrecta. & Se interpreta para comprender obstáculos y orientar revisión, manteniendo precisión. \\\hline
Diversidad & Ritmo, instrumento y condiciones uniformes como tendencia. & Reconocimiento de contextos, barreras y apoyos para sostener participación y aprendizaje. \\\hline
\end{longtable}}



{\small\setstretch{1.1}\renewcommand{\arraystretch}{1.25}
\begin{longtable}{|>{\raggedright\arraybackslash}p{\dimexpr0.3\linewidth-2\tabcolsep-2\arrayrulewidth\relax}|>{\raggedright\arraybackslash}p{\dimexpr0.35\linewidth-2\tabcolsep-2\arrayrulewidth\relax}|>{\raggedright\arraybackslash}p{\dimexpr0.35\linewidth-2\tabcolsep-2\arrayrulewidth\relax}|}\hline
\rowcolor{iiiepePale} \textbf{Criterio} & \textbf{Tendencia tradicional} & \textbf{Enfoque de la NEM} \\\hline
\endfirsthead
\hline
\rowcolor{iiiepePale} \textbf{Criterio} & \textbf{Tendencia tradicional} & \textbf{Enfoque de la NEM} \\\hline
\endhead
Relación con la comunidad & La actividad escolar puede quedar separada de su entorno. & La comunidad se considera referente de problemas, saberes y participación. \\\hline
Evidencia de comprensión & Énfasis en respuesta y reproducción. & Explicación, representación, transferencia y revisión del razonamiento. \\\hline
Responsabilidad & La evaluación se concentra en el docente. & Se promueve participación de estudiantes y diálogo, sin eliminar responsabilidad docente. \\\hline
Límite del contraste & Existen explicaciones directas útiles y prácticas tradicionales diversas. & Usar proyectos o tecnología no demuestra por sí mismo integración curricular ni aprendizaje. \\\hline
\end{longtable}}


\section{Aplicación a mi práctica}
Mi trayectoria en mantenimiento, mecatrónica e ingenierías ha priorizado organización del conocimiento, secuencias y resultados verificables. Estos recursos pueden contribuir a una enseñanza rigurosa; requieren revisarse cuando el control del procedimiento ocupa el lugar de la comprensión. Mi acercamiento humanista consiste en preguntar cómo interpreta el estudiante el problema y qué apoyo necesita para explicar una decisión.
En una actividad de calibración, conservaría la explicación directa de unidades y seguridad cuando sea necesaria. Después pediría formular hipótesis, comparar representaciones y revisar un dato discrepante. La evaluación integraría procedimiento y argumento, con una oportunidad de corrección. Así, la comparación orienta decisiones concretas y evita sustituir una práctica por una etiqueta.
\section{Una comparación que cambia la evaluación}
Considérese la ecuación $0.5x+2=0.2x+3.8$. Un ejercicio de respuesta final puede comprobar que el estudiante obtiene $x=6$, pero no revela si comprende equivalencia. Una tarea de desempeño añade tres solicitudes: explicar qué operación conserva la igualdad, comprobar en ambos miembros y construir un problema que produzca esa relación. La diferencia está en la evidencia disponible para decidir el siguiente apoyo.
La comparación de Scallon (2004) ayuda a distinguir producto, proceso y participación en la valoración. El marco de la SEP (2022a, 2022b) añade finalidades curriculares y comunitarias que no se reducen a cambiar el instrumento. Un examen breve puede aportar evidencia útil y un proyecto puede quedarse en reproducción: la calidad depende de qué razonamiento exige la tarea y cómo se utiliza lo observado.
En mi práctica técnica conservaría demostraciones claras cuando sean necesarias y las acompañaría de explicación del estudiante. El criterio humanista es mantener una expectativa común de comprensión con apoyos accesibles, evitando que velocidad, experiencia previa con herramientas o disponibilidad tecnológica se conviertan en sustitutos del aprendizaje matemático.

\section{Conclusión}
El cambio relevante consiste en organizar condiciones para que más estudiantes comprendan y participen con rigor. La NEM ofrece un marco para vincular conocimiento, comunidad y formación integral; su concreción requiere diseño, acompañamiento y evidencia. Mi criterio será revisar qué decisiones amplían comprensión y autonomía, manteniendo los estándares matemáticos y técnicos pertinentes a cada nivel.

\needspace{5\baselineskip}\section*{Referencias}
{\setstretch{1.15}\setlength{\parskip}{6pt}
\par\hangindent=1.27cm\hangafter=1  Scallon, G. (2004). \textit{L’évaluation des apprentissages dans une approche par compétences}. Éditions du Renouveau Pédagogique. Comparación de prácticas incluida en EVA.
\par\hangindent=1.27cm\hangafter=1  Secretaría de Educación Pública. (2022a). \textit{Plan de Estudio para la educación preescolar, primaria y secundaria}. SEP.
\par\hangindent=1.27cm\hangafter=1  Secretaría de Educación Pública. (2022b). \textit{Sugerencias metodológicas para el desarrollo de los proyectos educativos} [Extracto de \textit{Avance del contenido para el libro del docente. Primer grado}, material en proceso de edición, pp. 64--82, 23--24]. SEP.
}

\end{document}
